\chapter{Development Workflows}
\label{ch:workflows}

Chapter~\ref{ch:installing} noted, in passing, that Vim is Git's default editor: the program invoked automatically whenever Git needs a human to write or review something it cannot resolve on its own. This chapter takes that fact seriously and works through what it actually means in practice, since a substantial fraction of a working programmer's contact with Vim happens exactly here, often before they have deliberately chosen to learn the editor at all.

\section{Commit Messages}

When \texttt{git commit} is run without \texttt{-m}, Git opens the configured editor on a temporary file, pre-populated with a commented template, and expects a commit message written and saved back to that file, closing the editor to complete the commit or leaving the message empty (or the file unsaved, depending on configuration) to abort it. This file is an ordinary Vim buffer in every respect covered so far in this book: Insert mode types the message, Normal-mode commands from Parts II and III edit it, and \texttt{:wq} from Chapter~\ref{ch:installing} saves and exits to complete the commit. The lines beginning with \texttt{\#} are comments, stripped automatically by Git before the message is recorded, and are generally left alone rather than deleted, since they typically show which files are staged and provide useful confirmation of exactly what the commit is about to include.

\section{Interactive Rebase}

\texttt{git rebase -i} \textit{commit} opens a considerably more structured buffer: one line per commit being rebased, each beginning with an action word, \texttt{pick} by default, followed by that commit's abbreviated hash and message. This file is not merely a message to be written; it is a plan, and editing it edits what Git is about to do to the repository's history.

The most common edit, changing several \texttt{pick} lines to \texttt{squash} (folding a commit into the one before it) or its abbreviated form \texttt{s}, is an excellent worked example of the repeat command from Chapter~\ref{ch:language} applied to real, structured work rather than to prose. Position the cursor on the word \texttt{pick} on the first line to be squashed and press \texttt{cw}, the change-word command from Chapter~\ref{ch:operators}, type \texttt{squash}, and return to Normal mode. Move to the next \texttt{pick} line needing the same treatment and press \texttt{.}: because \texttt{cw} plus the typed replacement constitutes a single recorded change, per the compound-operator behavior of \texttt{c} noted in Chapter~\ref{ch:operators}, the period command replays the entire substitution, word and replacement together, against the new line. A rebase touching many commits becomes a short, rhythmic sequence: move to the next line, return to its first column with \texttt{0} from Chapter~\ref{ch:coordinates}, press \texttt{.}, and repeat, considerably faster and less error-prone than manually retyping \texttt{squash} on each line by hand.

Reordering commits is equally direct: \texttt{dd} removes the current line into a register exactly as it would in any other buffer, per Chapter~\ref{ch:operators}, and \texttt{p} or \texttt{P} places it below or above the cursor at its new position. Nothing about this is metaphorical or specific to rebase files; it is the same line-oriented deletion and placement covered throughout Part III, applied here to a file whose lines happen to represent commits rather than prose.

\section{Merge Conflicts}

When Git cannot automatically reconcile two branches' changes to the same region of a file, it leaves that file with the conflict marked directly in its text, delimited by \texttt{<<<<<<<}, \texttt{=======}, and \texttt{>>>>>>>}. Opening such a file in Vim presents the conflict as exactly what it is: ordinary text, with unusual but perfectly regular structural markers, resolvable with the same commands covered throughout this book rather than any conflict-specific vocabulary. A search for \texttt{<<<<<<<} with \texttt{/} from Chapter~\ref{ch:coordinates} navigates directly to the next conflict; the operators and text objects of Parts II and III delete whichever side of the conflict is being discarded, along with the marker lines themselves, leaving a single, coherent resolution behind.

A pattern worth knowing specifically for this task: \texttt{:g/\textasciicircum<<<<<<</,/\textasciicircum=======/d} deletes every line from a conflict's opening marker through its middle marker, discarding the current branch's version of a conflicted section in a single command, if the incoming branch's version is the one to keep; the corresponding deletion from the middle marker through the closing one discards the incoming version instead. Constructing this command is a direct application of the pattern-defined-range technique from Chapter~\ref{ch:ex}, combined with the deletion command from \texttt{:g}, exactly as covered in Chapter~\ref{ch:global}, and is considerably faster than manually selecting and deleting each side of a conflict by hand when a file contains many conflicted sections following the same resolution pattern.

\section{Comparing Files with \texttt{vimdiff}}

\texttt{vimdiff} \textit{file1} \textit{file2} (or, equivalently, \texttt{vim -d}) opens two or more files side by side in vertically split windows, per Chapter~\ref{ch:buffers}, with differing regions highlighted directly in the buffer text. \texttt{]c} and \texttt{[c} jump to the next and previous difference; \texttt{do} (diff obtain) pulls the other window's version of the current difference into the current window, and \texttt{dp} (diff put) pushes the current window's version to the other. Git can be configured to use \texttt{vimdiff} as its merge tool directly (\texttt{git config merge.tool vimdiff}), at which point \texttt{git mergetool} opens conflicts in this side-by-side view rather than as inline markers within a single buffer, a genuinely different and, for some readers, more legible way to work through the same underlying conflict-resolution task covered above.

\section{The Quickfix List and \texttt{make}}

\texttt{:make} runs the project's configured build command (by default, simply \texttt{make}, though the \texttt{makeprg} option can point it at any other build tool) and captures its output into the quickfix list: a structured list of file-and-line-number references, generally corresponding to compiler errors or warnings. \texttt{:copen} opens the quickfix list in its own window, one entry per line, and pressing return on any entry jumps directly to the corresponding file and line; \texttt{:cnext} and \texttt{:cprevious} step through the list without needing to return to the quickfix window each time. This transforms the otherwise manual work of reading a build's error output and separately navigating to each reported location into a single, structured workflow, staying entirely within Vim rather than switching between a terminal's scrollback and the editor.

\section{Searching a Project with \texttt{grep}}

\texttt{:grep} \textit{pattern} runs an external \texttt{grep} (or, on some systems, a faster alternative it has been configured to use instead) across the project and populates the quickfix list with every matching line, exactly as \texttt{:make} does with build errors; the same \texttt{:copen}, \texttt{:cnext}, and \texttt{:cprevious} commands then navigate the results. A closely related built-in alternative, \texttt{:vimgrep} \textit{pattern} \texttt{**/*.py}, searches using Vim's own regex engine rather than shelling out to an external program, which trades some speed on very large codebases for exact consistency with the regex syntax already covered in Chapter~\ref{ch:substitution}, since a pattern that works in \texttt{:s} is guaranteed to work identically in \texttt{:vimgrep}, which is not always true of an external \texttt{grep}'s own regex dialect.

\section{Editing History as Editing Text}

The material in this chapter has one idea running underneath all of it, worth stating explicitly now that the individual workflows have been covered: a commit message, a rebase plan, a conflicted file, and a build's error output are all, structurally, just text, and every one of them yields to the same grammar covered throughout this book rather than requiring any Git-specific or build-specific vocabulary of its own. The apparent gap between "editing a document" and "managing a project's history" turns out not to be a gap in kind at all, only in what the text being edited happens to represent. A reader fluent in the material of Parts II through V is already fluent in everything this chapter has covered; what this chapter has added is only the recognition of where that fluency applies.
